\appendixpart{app:grid}
\setcounter{table}{0}
\renewcommand{\thetable}{А.\arabic{table}}
\begin{center}\textbf{Результаты параметризации}\end{center}

Все значения отклонений приведены в процентах. В последнем столбце дан максимум по классу графов. Для каждого графа и сочетания параметров проведено по 20 запусков.

\begin{table}[H]\centering
\caption{Результаты для всех комбинаций параметров}
\label{tab:grid}
{\small\setstretch{1.05}\setlength{\tabcolsep}{6pt}
\begin{tabular}{@{}rrr rrrrrrr@{}}
\toprule
 &&& \multicolumn{2}{c}{$G_1$} & \multicolumn{2}{c}{$G_2$} & \multicolumn{2}{c}{$G_3$} & Класс \\
\cmidrule(lr){4-5}\cmidrule(lr){6-7}\cmidrule(lr){8-9}
$\alpha$ & $\beta$ & $\rho$ &
$\delta_{\max}$ & $\widetilde\delta$ &
$\delta_{\max}$ & $\widetilde\delta$ &
$\delta_{\max}$ & $\widetilde\delta$ & $\delta_{\max}$ \\
\midrule
\tablerows{results/parameters.tex}
\bottomrule
\end{tabular}}
\end{table}

\appendixpart{app:graphs}
\setcounter{lstlisting}{0}
\renewcommand{\thelstlisting}{Б.\arabic{lstlisting}}
\begin{center}\textbf{Исходные данные параметризации}\end{center}

Первая строка каждого файла содержит число портов и бюджет, вторая --- имена портов. Последующие строки содержат пары «время, стоимость» для $i<j$, в порядке возрастания $i$, затем $j$.

\lstinputlisting[language={},numbers=none,xleftmargin=0mm,linewidth=\linewidth,
basicstyle=\ttfamily\footnotesize\setstretch{1.15},
caption={Исходные данные графа $G_1$}]{data/G1.txt}
\lstinputlisting[language={},numbers=none,xleftmargin=0mm,linewidth=\linewidth,
basicstyle=\ttfamily\footnotesize\setstretch{1.15},
caption={Исходные данные графа $G_2$}]{data/G2.txt}
\lstinputlisting[language={},numbers=none,xleftmargin=0mm,linewidth=\linewidth,
basicstyle=\ttfamily\footnotesize\setstretch{1.15},
caption={Исходные данные графа $G_3$}]{data/G3.txt}
